\documentclass[10pt,fleqn,a4paper]{jsarticle}
\usepackage[dvipdfmx]{graphicx,color}
\usepackage{ascmac,amsmath,amssymb,amstext}
\usepackage{tikz,multicol,float,tikz-3dplot}
\usetikzlibrary{positioning,intersections,calc,arrows.meta,math,angles,patterns}
\usepackage{zogeny,ceo,comment}
\newcommand{\theme}[1]{【\text{#1}】}
\newcommand{\ans}[1]{\mbox{\boldmath{$#1$}}}
\renewcommand{\baselinestretch}{1.3} 
\newcommand{\hokusomu}[1]{
\begin{tikzpicture}[scale=1]
\filldraw[fill=black,draw=black](0,0)--({1-0.2},{0})--({1-0.2},{1-0.2})--(0,{1-0.2})--(0,0);
\draw[-,thick](0,0)--({1-0.2},{0})--({1/2-0.2},{1-0.2})--(0,{1-0.2})--(0,0);
\node[white,font=\Huge]at({1/2-0.1},{1/2-0.1})(label){$\mathbf{#1}$};
\end{tikzpicture}}
\newcommand{\hokukai}{
    \begin{tikzpicture}
  \draw[line width=0.2mm](-0.73,{-1/4+0.04})--(-0.73,{1/4-0.03})--(0.73,{1/4-0.03})--(0.73,{-1/4+0.04})--cycle;
  \draw[line width=0.6mm](-0.73,{-1/4+0.03})--(-0.73,{1/4-0.02});
  \draw[line width=0.6mm](0.73,{-1/4+0.03})--(0.73,{1/4-0.02});
  \node at(0,0){\kurosankakub \ans{解答}\kurosankakua};
\end{tikzpicture}
} %全国大学入試問題の「解答」記号
\tdplotsetmaincoords{60}{110}%{xy平面からどれだけ上にいるか(90度から-90度)}{z軸中心にどれだけ左右に回転するか}
\title{数学の核心-理系}
\author{大島　遙斗}
\begin{document}
\maketitle
\newpage
\tableofcontents
\newpage
\noindent{\Large \ans{はじめに}}\\
このテキストでは,数学3の内容を扱います.特に,「収束・発散速度」「接線近似」「力学系」「偏微分」等々を扱います.\\
頑張りましょう.大島
\newpage
\section{指数対数関数の微分}

\begin{itembox}
[l]{\enpitu 関数の極限}
$c$を定数とする.変数$x$を定数$a$にどのような近づけ方をしても,
$|f(x)-c|$の値が限りなく0に近づくとき,
\begin{center}
$\displaystyle \lim_{x\to a}f(x)=c$
\end{center}
と表す.特に,
\begin{center}
$
\displaystyle
\lim_{x\to a}f(x)=c
\Longleftrightarrow
\begin{cases}
\displaystyle\lim_{x\to a+0}f(x)=c\\[2mm]
\displaystyle\lim_{x\to a-0}f(x)=c
\end{cases}
$
\end{center}
である.
\end{itembox}

\newpage

\noindent\ba{1$\bullet$1}\
\par\noindent\kakkoichi 2は3の何乗ですか?\
\par\noindent\kakkoni 自然対数の底$e$とは何ですか?\
\par\noindent\kakkosan 次の極限の公式がすべて同じものに見えますか?\
\begin{align*}
&\lim_{x\to0}\dfrac{e^x-1}{x}=1,
&&\lim_{x\to0}\dfrac{\log(x+1)}{x}=1,\\
&\lim_{x\to0}(1+x)^{\frac{1}{x}}=e,
&&\lim_{x\to\infty}\left(1+\dfrac{1}{x}\right)^x=e,
&&\lim_{x\to-\infty}\left(1+\dfrac{1}{x}\right)^x=e.
\end{align*}
\par\noindent\kakkoshi $(\log x)'=\dfrac{1}{x}$が当たり前に見えますか?\\

\noindent\ba{1$\bullet$2}\
次の極限を求めよ.\
\par\noindent\kakkoichi $\displaystyle\lim_{\theta\to0}\dfrac{\sin2\theta}{\tan3\theta}$\
\par\noindent\kakkoni $\displaystyle\lim_{x\to0}\dfrac{3^x-1}{x}$\
\par\noindent\kakkosan $\displaystyle\lim_{x\to0}\dfrac{\log(1-\sin x)}{x}$\
\par\noindent\kakkoshi $\displaystyle\lim_{x\to0}\dfrac{|x|}{x}$\
\par\noindent\kakkogo $\displaystyle\lim_{x\to0}(1-x)^{-\frac{1}{x}}$\
\par\noindent\kakkoroku $\displaystyle\lim_{x\to0}(1+x)^{\frac{1}{x^2}}$\

\newpage

\section{発散速度と関数のグラフ}

\begin{itembox}
[l]{\enpitu 不等式と極限}
\noindent\kakkoichi
$
\begin{cases}
十分大きなxに対し,g(x)<f(x)<h(x)\\
\displaystyle\lim_{x\to\infty}g(x)=\lim_{x\to\infty}h(x)=\alpha\quad(\alpha\text{は定数})
\end{cases}
$
ならば,
\begin{center}
$\displaystyle\lim_{x\to\infty}f(x)=\alpha$
\end{center}
である.\\
\noindent\kakkoni
$
\begin{cases}
十分大きなxに対し,g(x)<f(x)\\
\displaystyle\lim_{x\to\infty}g(x)=\infty
\end{cases}
$
ならば,
\begin{center}
$\displaystyle\lim_{x\to\infty}f(x)=\infty$
\end{center}
である.
\end{itembox}

\begin{itembox}
[l]{\enpitu 発散速度}
\begin{center}
$
\begin{aligned}
\cdots&\ll\log x\ll\cdots\ll\sqrt{x}\ll x\ll x^2\ll x^3\ll\cdots\\
&\ll2^x\ll e^x\ll3^x\ll\cdots\ll x^x\ll\cdots
\end{aligned}
$
\end{center}
\end{itembox}

\newpage

\noindent\ba{2$\bullet$1}\
\par\noindent\kakkoichi $e^x\geq1+x$を示せ.\
\par\noindent\kakkoni $\displaystyle\lim_{x\to\infty}\dfrac{e^x}{x}=\infty$を示せ.\
\par\noindent\kakkosan 1より大きい任意の定数$a$と正の任意の定数$b$に対して,
$\displaystyle\lim_{x\to\infty}\dfrac{a^x}{x^b}=\infty$を示せ.\\

\noindent\ba{2$\bullet$2}\
$f(x)=x^x\ (x>0)$とする.増減,凹凸を調べ,$y=f(x)$のグラフを描け.\\

\noindent\ba{2$\bullet$3}\
$a$は1より大きい定数とし,$xy$平面上の点$(a,0)$をA,点$(a,\log a)$をB,
曲線$y=\log x$と$x$軸の交点をCとする.さらに$x$軸,線分BAおよび曲線$y=\log x$で
囲まれた部分の面積を$S_1$とする.\
\par\noindent\kakkoichi $1\leq b\leq a$となる$b$に対し,点$(b,\log b)$をDとする.
四辺形ABDCの面積が$S_1$にもっとも近くなるような$b$の値と,
そのときの四辺形ABDCの面積$S_2$を求めよ.\
\par\noindent\kakkoni $a\to\infty$のときの$\dfrac{S_2}{S_1}$の極限値を求めよ.\
\par\noindent\hfill（東京大）\

\newpage

\section{「収束速度」のはなし}

\begin{itembox}
[l]{\enpitu ランダウの記号}
\begin{align*}
\lim_{x\to a}f(x)=0,\qquad
\lim_{x\to a}g(x)=0,\qquad
\lim_{x\to a}\dfrac{f(x)}{g(x)}=0
\end{align*}
のとき,「$x=a$の近くで$f(x)$は$g(x)$より高位の無限小である」といい,
\begin{center}
$f(x)=o(g(x))\quad(x\to a)$
\end{center}
と書く.
\end{itembox}

\begin{itembox}
[l]{\enpitu 微分の定義}
関数$y=f(x)$に対し,
\begin{align*}
f'(x)
&=\lim_{\Delta x\to0}\dfrac{f(x+\Delta x)-f(x)}{\Delta x}
=\lim_{\Delta x\to0}\dfrac{\Delta y}{\Delta x}\\
&\Longleftrightarrow \lim_{\Delta x\to0}\dfrac{\Delta y-f'(x)\Delta x}{\Delta x}=0\\
&\Longleftrightarrow \Delta y-f'(x)\Delta x=o(\Delta x)\quad(\Delta x\to0)\\
&\Longleftrightarrow \Delta y=f'(x)\Delta x+o(\Delta x)\quad(\Delta x\to0)\\
&\Longleftrightarrow dy=f'(x)dx.
\end{align*}
\end{itembox}

\begin{itembox}
[l]{\enpitu 微積分学の基本定理}
$a$が定数のとき,
\begin{center}
$\displaystyle\dfrac{d}{dx}\int_a^x f(t)\,dt=f(x)$
\end{center}
である.
\end{itembox}

\newpage

\noindent\ba{3$\bullet$1}\
関数$y=f(x)=x^3$について,$x$の増分$\Delta x$に対する$y$の増分を$\Delta y$とするとき,\
\par\noindent\kakkoichi $\Delta y$を$\Delta x$で表せ.\
\par\noindent\kakkoni $dy$を$dx$で表せ.\\

\noindent\ba{3$\bullet$2}\
$a$を定数とするとき,微積分学の基本定理
\begin{center}
$\displaystyle\dfrac{d}{dx}\int_a^x f(t)\,dt=f(x)$
\end{center}
を説明せよ.\\

\noindent\ba{3$\bullet$3}\
積の微分公式
\begin{center}
$\displaystyle\dfrac{d}{dx}\{f(x)g(x)\}=f'(x)g(x)+f(x)g'(x)$
\end{center}
を証明せよ.\\

\noindent\ba{3$\bullet$4}\
\par\noindent\kakkoichi 半径$r$の円の面積$\pi r^2$を$r$で微分すると,
その円の周長$2\pi r$になるのはなぜか.\
\par\noindent\kakkoni 半径$r$の球の体積$\dfrac{4}{3}\pi r^3$を$r$で微分すると,
その球の表面積$4\pi r^2$になるのはなぜか.\

\end{document}